\documentclass[10pt]{article}
\usepackage[margin=20mm]{geometry}
\usepackage{amsmath,amssymb,mathtools,fontspec,unicode-math,xurl,hyperref}
\IfFontExistsTF{DejaVu Serif}{\setmainfont{DejaVu Serif}}{\setmainfont{Latin Modern Roman}}
\setmathfont{latinmodern-math.otf}
\setlength{\parindent}{0pt}\setlength{\parskip}{6pt}\sloppy
\title{Initial review of
On the \(p\)-adic weight-monodromy conjecture
for complete intersections in toric varieties}
\author{}
\date{26 October 2026}
\begin{document}
\maketitle
Abstract: This report concerns Federico Binda, Hiroki Kato and Alberto Vezzani’s 36-page arXiv:2207.00369v2 (10 June 2025). It establishes a false unqualified connectedness assertion and explains precisely the consequent scope problem for the one-dimensional top-cohomology step. It does not disprove weight–monodromy. A second, local issue is the omitted unit summand in two torus-motive identities. The connected case of the main splitting argument survives these checks.\par
\section*{Identity, outcome and review limits}
The assigned PDF and author TeX are bound to arXiv:2207.00369v2 by the supplied source manifest [paper]. The associated journal DOI is 10.1007/s00222-025-01344-x. A publisher version of record was not acquired; no issue below is asserted to occur in unchanged journal-final text. All page numbers below are one-based PDF pages, identical to printed page numbers. The author-source line numbers are one-based.\par
The outcome is Major: proof scope requires correction, with one additional Minor local proof identity. The Major assessment concerns the unrestricted scope represented by Theorem 1.2 and Remark 5.8, not a counterexample to the weight–monodromy theorem. If the word variety in Theorem 5.6 is intended to impose geometric integrality, its connected-case argument is outside this objection; the explicit broader connectedness claim still fails.\par
The entire extracted 2,121-line PDF text, including the appendix and references, was read consecutively. Fifteen original PDF pages were actually viewed: 1, 2, 8, 18–20, 22, 24–29, 31 and 32. The 2,530-line author TeX was parsed in full for a static lexical audit; the preamble and relevant matched source passages were inspected. This is not a claim that every author-source line was separately read or that the source was executed.\par
\section*{Claims, hypotheses and proof dependency map}
Theorem 1.2 concerns a smooth scheme-theoretic complete intersection in a projective smooth toric variety over a finite extension of \(\mathbb Q_p\). Theorem 5.6 states the geometric version over \(\mathbb C_p\). Quasi-purity of weight \(i\) means that the \(j\)-th monodromy-graded piece of \(H^i_{\rm HK}\) is Frobenius-pure of weight \(i+j\); the definition uses the appropriate linear power of semilinear Frobenius. These conventions matter when transporting monodromy across ramified extensions.\par
The main chain inspected is:\par
1. Theorem 2.20 identifies log motives with the motives of their log special fibers. Proposition 2.29 and Corollary 2.33 connect this with tilting and rigid generic fibers.\par
2. Proposition 3.8 constructs the geometric Hyodo–Kato realization. Corollary 3.13 identifies the tilted realization on semistable models, with monodromy normalized by the nonzero ramification index.\par
3. Proposition 4.1 uses continuity and compactness of the motive of a smooth quasi-compact, quasi-separated family to obtain a small neighborhood with the same motive. Corollary 4.6 applies this to a smooth scheme-theoretic complete intersection.\par
4. Scholze’s approximation supplies an irreducible equal-characteristic subvariety \(Z\) of the same dimension inside the inverse image of that neighborhood. An alteration \(Z'\) and Construction 5.5 produce \[\alpha:R\Gamma_{\rm HK}(Y)\longrightarrow
 R\Gamma_{\rm HK}^{\flat}(Z').\]\par
5. A positive cycle gives nonzero top-degree pullback. If both top-degree groups are one-dimensional, this is an isomorphism. Poincare duality then splits every \(\alpha^i\), and Proposition 5.3 supplies quasi-purity on the target.\par
The passage from nonzero to isomorphism is the precise point affected by Problem 1. Section 3.3 is expressly optional for this main proof (page 20); its arithmetic realization is distinct from the geometric construction used in Section 5.\par
\section*{Problem 1: connectedness and the scope of top-degree splitting}
\subsection*{Location and exact assertion}
Remark 5.8, page 29, asserts geometric connectedness for every positive-dimensional complete intersection in a smooth projective toric variety, citing SGA 2, XII.3.5 and Cohen–Macaulayness. Theorem 5.6, page 28, uses one-dimensionality of both sides of \(\alpha^{2d}\). The relevant author-source lines are 1740–1768 and 1776–1778; Theorem 1.2 is at lines 431–434.\par
The geometric setup immediately before Corollary 4.6 (page 26; source lines 1597–1617) allows sections of general line bundles \(\mathcal O(D_i)\). There is no ampleness condition on those bundles. Thus a projective ambient toric variety cannot silently supply positivity of every defining divisor.\par
\subsection*{A complete counterexample to the connectedness assertion}
Take \[K=\mathbb Q_3,\qquad L=K(\sqrt2),\qquad
 X=\mathbb P^1_K\times\mathbb P^1_K,\qquad
 Y=V(x_0^2-2x_1^2).\] The equation is a global section of \(\mathcal O_X(2,0)\). The polynomial \(t^2-2\) is irreducible modulo \(3\), so \(L/K\) is a separable quadratic extension and the zero scheme in the first factor is \(\operatorname{Spec}L\). Consequently \[Y\simeq\mathbb P^1_L
 \quad\hbox{as a $K$-scheme}.\] This verifies all the relevant conditions directly: \(X\) is smooth, projective and toric; \(Y\) has dimension one and codimension one; its single equation is a nonzero divisor, so it is a scheme-theoretic complete intersection. It is integral over \(K\). It is smooth over \(K\), since \(\operatorname{Spec}L\to
\operatorname{Spec}K\) is finite etale and projective space is smooth.\par
Over \(\mathbb C_3\), the two distinct roots give \[Y_{\mathbb C_3}
   =\{\sqrt2\}\times\mathbb P^1_{\mathbb C_3}
     \ \amalg\
     \{-\sqrt2\}\times\mathbb P^1_{\mathbb C_3}.\] In local coordinates the two defining linear factors have nonzero derivative; there is no nilpotent or singular intersection hidden in the example. The two components are disjoint and open and closed. Thus \(Y\) is not geometrically connected, despite being integral over \(K\) and of positive dimension.\par
The same issue can be exhibited directly over an algebraically closed field by \(V(x_0x_1)\subset\mathbb P^1\times\mathbb P^1\). The first example also excludes the answer that ordinary integrality over the original finite extension of \(\mathbb Q_p\) suffices.\par
\subsection*{Failure of the cited inference}
The restriction of \(\mathcal O_X(2,0)\) to any second-factor fiber has degree zero, so this line bundle is not ample. The cited SGA 2, XII.3.5 explicitly assumes an ample invertible sheaf [sga]. Cohen–Macaulayness addresses a depth hypothesis; it does not imply this positivity hypothesis. The primary text of that corollary was checked in the 2005 reedition, PDF page 128, printed page 120.\par
Write \(F=\mathbb Q_3^{\rm nr}\), the coefficient field in the paper’s geometric realization. Additivity and the projective-line calculation give \[\dim_F H^0_{\rm HK}(Y_{\mathbb C_3})=2,\qquad
 \dim_F H^2_{\rm HK}(Y_{\mathbb C_3})=2,\qquad H^1_{\rm HK}=0.\] For the single smooth irreducible proper curve \(Z'\) selected in the proof, the top group has dimension one. A nonzero linear map \(F^2\to F\) cannot be an isomorphism. Its rank is one and its kernel is nonzero. Likewise, no split injection from this \(H^0(Y)\) into \(H^0(Z')\) can exist. Poincare duality cannot remove that rank obstruction. This proves failure of the stated single-component inference whenever the broader, possibly geometrically disconnected scope is intended.\par
\subsection*{Counterevidence, impact and bounded repairs}
This example satisfies weight–monodromy: each geometric component is a projective line, \(N=0\), and the cohomology has the expected weights. It is not a counterexample to the conclusion. For geometrically connected smooth proper \(Y\), the top group is one-dimensional, and the displayed nonzero-map and duality argument has no defect from this test. If variety over \(\mathbb C_p\) means integral, Theorem 5.6 already belongs to that case. The unqualified Remark 5.8 and the passage from the broader statement in the introduction still require explicit treatment.\par
A safe scope repair is to impose geometric connectedness in the main statements and delete or appropriately restrict Remark 5.8. This narrows the advertised result. A positivity-based repair would require the actual ample-divisor hypotheses needed for an iterated Lefschetz argument; ambient projectivity alone is insufficient.\par
To retain all smooth complete intersections, a separate argument must handle every geometric component and construct enough comparison targets to obtain a splitting on the full direct sum. Simply working componentwise is not a complete repair: a component of a global complete intersection need not inherit the same given global complete-intersection presentation. The current review does not supply that missing general approximation argument. The Major classification reflects this unproved scope extension, not failure of the connected theorem.\par
\section*{Problem 2: the full torus motive loses its unit summand}
\subsection*{Location and check}
In the proof of Proposition 3.29, page 22 (source line 1411), two equalities replace the motive of a product with \(\mathbb T^k\), respectively \(\mathbb T^1\), by a single twisted and shifted motive. The definition on page 8 identifies \(\mathbb T^1\) with the actual rigid analytic torus associated with the formal multiplicative group; it is not a notation for a reduced motive.\par
There is a convention-independent obstruction. The unit section \(1:\operatorname{Spa}L\to\mathbb T^1\) and structural projection split off the unit motive. Thus, writing \(\widetilde M(\mathbb T^1)\) for the complementary reduced motive, \[M(X\times\mathbb T^1)
 \simeq M(X)\oplus
       \bigl(M(X)\otimes\widetilde M(\mathbb T^1)\bigr).\] With the usual homological Tate convention the second summand is \(M(X)(1)[1]\). The unit summand is required independently of how the Tate twist is normalized.\par
For \(X=\operatorname{Spa}L\), overconvergent de Rham cohomology of the torus has a constant class in degree zero and the class \(d u/u\) in degree one. The latter is not exact: termwise differentiation of a Laurent series has zero coefficient of \(u^{-1}\,du\). For nonzero exponents, division by the exponent preserves overconvergence after decreasing the radius, yielding the usual one-dimensional degree-one quotient. Hence the total cohomology dimension is two. A single shifted Tate object has total dimension one. This disproves the displayed full-motive identity, not merely a choice of sign or grading.\par
\subsection*{Counterevidence, impact and repair}
The purpose of both identities is only to establish membership in \(\operatorname{RigDA}_{\rm gr}(L)\). This category is closed under sums, shifts and twists by Definition 3.22. Replacing the false identity by the full direct-sum decomposition still establishes the needed membership, including the Mayer–Vietoris term. For \(k\) tori, use the tensor product of \(k\) split two-term motives, not a single summand.\par
One could rescue the formula by explicitly using a reduced torus motive, but the displayed expression \(M(X\times\mathbb T^k)\) and the unreduced Mayer–Vietoris cover do not make that replacement. The flaw is therefore local and repairable. It does not establish failure of Proposition 3.29, and Section 3.3 is not needed for the main theorem. The appropriate severity is Minor.\par
\section*{Other checks and unresolved dependencies}
The normalization \(N\mapsto N/e\) in Corollary 3.13 was checked: for \(c\ne0\), kernels and images of \((cN)^j\) equal those of \(N^j\), so the monodromy filtration is unchanged. Multiplication by \(c^j\) preserves the required graded isomorphisms. This resolves a possible normalization objection.\par
For Corollary 4.6, smoothness of the fiber and a codimension-\(c\) regular presentation imply independence of the \(c\) differentials along the fiber. The map to the \(c\)-dimensional polydisc is therefore smooth near it. Quasi-compactness and a finite trivializing cover permit one sufficiently small common radius. The distinction between scheme-theoretic and merely set-theoretic intersection is essential to this Jacobian argument.\par
In the connected case, the splitting in Theorem 5.6 is also directly verified algebraically. Normalize the target trace by the nonzero top-degree scalar. For the adjoint \(r\) of \(\alpha^{2d-i}\), multiplicativity of \(\alpha\) gives \(\langle r\alpha x,y\rangle=\langle x,y\rangle\) for all \(x,y\). Nondegeneracy yields \(r\alpha=1\), compatibly with \((\varphi,N)\). A direct summand inherits the monodromy filtration and quasi-purity.\par
The fine-log reduction, approximation of non-fine charts, rigid-etale descent and all higher-categorical coherence results were traced through the internal dependency chain, but the cited long proofs in AGV22, Ayoub and the comparison literature were not independently re-proved. The equal-characteristic theorem of Lazda–Pal and alteration inputs remain external dependencies. In the appendix, the stated construction on weight-graded cohomology was distinguished from a full complex-level monodromy enhancement referred to BGV23. No additional confirmed defect is assigned to those unverified external portions.\par
\section*{Supporting checks}
These are bounded sanity checks; the counterexample and general monodromy argument above are analytic proofs, not conclusions from numerical sampling.\par
The static source scan recorded packages, the `xy` input, and all matched connectedness, ampleness and torus-motive contexts. No direct execution primitive matched the documented lexical patterns. This limited scan is not a security proof. Absence of author computation code is not treated as a defect.\par
The final report is compiled from its own exact current TeX with the authorized Tectonic 0.17.0. The build record binds pre/post hashes of the source, compiler and complete task-local dependency-cache inventory, command lines, exit status, logs and resulting PDF. The separate page-check record identifies each actually viewed final report page and its PNG hash. Coverage, findings and provenance are bound by hashes in the per-paper delivery state. Delivery completion is not scientific acceptance: `human\_reviews=[]`, `scientific\_decision=pending`.\par
9 F. Binda, H. Kato and A. Vezzani, On the \(p\)-adic weight-monodromy conjecture for complete intersections in toric varieties, arXiv:2207.00369v2, 10 June 2025. https://arxiv.org/abs/2207.00369. Assigned PDF SHA-256: `150ba49710056b7de62ca91467cd7e4cde2aaba6a053a5f83bd0c35e0addfea9`. A. Grothendieck et al., Cohomologie locale des faisceaux coherents et theoremes de Lefschetz locaux et globaux (SGA 2), 2005 reedition, XII, Corollary 3.5, printed p. 120. https://arxiv.org/pdf/math/0511279. The ample-sheaf hypothesis was checked in the primary text. P. Scholze, Perfectoid spaces, author-hosted preprint dated 19 November 2011, Corollary 8.8 and Theorems 1.14 and 9.6. https://people.mpim-bonn.mpg.de/scholze/PerfectoidSpaces.pdf. The stated weight–monodromy theorem explicitly assumes geometric connectedness; the approximation corollary permits general hypersurfaces.\par
\end{document}
